\documentclass{article}
\usepackage{amsmath}
\usepackage{amssymb}
\usepackage{geometry}
\usepackage{graphicx}
\usepackage{xepersian}
\settextfont{XB Niloofar}

\geometry{a4paper, margin=1in}

\title{\lr{Naive Bayes Classification} \\ (دسته‌بندی نایو بیز)}
\author{}
\date{}
\begin{figure}[h!]
  \centering
  \includegraphics[width=0.75\textwidth]{img_page_4_1.png}
  \caption{نمودار احتمال شرطی کلاس برای کلاس‌های مختلف}
\end{figure}

\begin{figure}[h!]
  \centering
  \includegraphics[width=0.75\textwidth]{img_page_5_1.png}
  \caption{نمودار احتمال پسین و تعیین کلاس نهایی}
\end{figure}
\begin{figure}[h!]
  \centering
  \includegraphics[width=0.75\textwidth]{img_page_5_2.png}
  \caption{تعیین مرز تصمیم بر اساس توابع چگالی احتمال}
\end{figure}

\begin{figure}[h!]
  \centering
  \includegraphics[width=0.75\textwidth]{img_page_6_1.png}
  \caption{توزیع احتمالات ویژگی‌ها در ویژگی‌های مختلف}
\end{figure}

\begin{figure}[h!]
  \centering
  \includegraphics[width=0.75\textwidth]{img_page_7_1.png}
  \caption{نمایش فرضیه ساده‌لوحانه استقلال ویژگی‌ها}
\end{figure}

\begin{figure}[h!]
  \centering
  \includegraphics[width=0.75\textwidth]{img_page_8_1.png}
  \caption{مرز تصمیم‌گیری غیرخطی طبقه‌بند بیز ساده روی داده‌های شبیه‌سازی‌شده}
\end{figure}

\begin{figure}[h!]
  \centering
  \includegraphics[width=0.75\textwidth]{img_page_12_1.png}
  \caption{توزیع احتمال گوسی برای متغیرهای پیوسته در دسته‌بند بیز}
\end{figure}
\begin{figure}[h!]
  \centering
  \includegraphics[width=0.75\textwidth]{img_page_12_2.png}
  \caption{توزیع احتمال چندجمله‌ای برای متغیرهای گسسته}
\end{figure}

\begin{figure}[h!]
  \centering
  \includegraphics[width=0.75\textwidth]{img_page_14_1.png}
  \caption{ساختار گرافیکی شبکه‌ای ساده‌لوحانه بیز (Naive Bayes Graph)}
\end{figure}

\begin{figure}[h!]
  \centering
  \includegraphics[width=0.75\textwidth]{img_page_27_1.png}
  \caption{ارزیابی و بررسی صحت طبقه‌بند بر روی داده‌های آزمون}
\end{figure}

\begin{figure}[h!]
  \centering
  \includegraphics[width=0.75\textwidth]{img_page_29_1.png}
  \caption{مزایا و محدودیت‌های الگوریتم بیز ساده}
\end{figure}






















\begin{document}

\maketitle

\section*{فهرست مطالب}
\begin{enumerate}
    \item مقدمه
    \item قضیه بیز (\lr{Bayes' Theorem})
    \item دسته‌بندی‌کننده‌های بیز (\lr{Bayes Classifiers})
    \item نایو بیز با ویژگی‌های متعدد (\lr{Naive Bayes with Multiple Features})
    \item مدل‌های زایشی (\lr{Generative Models})
    \item نایو بیز گاوسی (\lr{Gaussian Naive Bayes})
    \item نایو بیز چندجمله‌ای (\lr{Multinomial Naive Bayes})
    \item چه زمانی از نایو بیز استفاده کنیم؟
\end{enumerate}

\newpage

\section*{۱. مقدمه}
مدل‌های نایو بیز (\lr{Naive Bayes}) گروهی از الگوریتم‌های دسته‌بندی فوق‌العاده سریع و ساده هستند که اغلب برای مجموعه‌داده‌های با ابعاد بسیار بالا مناسب می‌باشند. به دلیل سرعت بالا و پارامترهای قابل تنظیم بسیار کم، این مدل‌ها به عنوان یک مدل پایه (\lr{baseline}) ساده و سریع برای مسائل دسته‌بندی بسیار مفید هستند.

\section*{۲. قضیه بیز}
قضیه بیز (\lr{Bayes' Theorem}) یک اصل اساسی در تئوری احتمال است که نحوه به‌روزرسانی باورهای ما درباره یک رویداد را بر اساس شواهد جدید توصیف می‌کند. فرمول ریاضی قضیه بیز به صورت زیر بیان می‌شود:
\[
P(A|B) = \frac{P(B|A) \times P(A)}{P(B)}
\]
که در آن:
\begin{itemize}
    \item $P(A|B)$ احتمال پسین (\lr{posterior probability}) است: احتمال وقوع رویداد A به شرط برقرار بودن B.
    \item $P(B|A)$ درست‌نمایی (\lr{likelihood}) است: احتمال وقوع رویداد B به شرط برقرار بودن A.
    \item $P(A)$ احتمال پیشین (\lr{prior probability}) است: احتمال اولیه وقوع رویداد A.
    \item $P(B)$ احتمال حاشیه‌ای (\lr{marginal probability}) است: احتمال کل وقوع رویداد B.
\end{itemize}
\begin{figure}[h!]
  \centering
  \includegraphics[width=0.75\textwidth]{img_page_2_1.jpeg}
  \caption{مدل مفهومی طبقه‌بندی ایمیل‌های اسپم و غیر اسپم بر اساس کلمات}
\end{figure}

\begin{figure}[h!]
  \centering
  \includegraphics[width=0.75\textwidth]{img_page_9_1.png}
  \caption{طبقه‌بندی بیز با ویژگی‌های چندگانه گسسته و پیوسته}
\end{figure}

\begin{figure}[h!]
  \centering
  \includegraphics[width=0.75\textwidth]{img_page_22_1.png}
  \caption{توزیع‌های نرمال چندمتغیره در طبقه‌بند بیز گوسی}
\end{figure}
\begin{figure}[h!]
  \centering
  \includegraphics[width=0.75\textwidth]{img_page_22_2.png}
  \caption{مرزهای تصمیم حاصل از طبقه‌بند بیز گوسی}
\end{figure}









\subsection*{مثال با تاس}
فرض کنید دو تاس داریم: یکی معمولی (با احتمال برابر برای آوردن اعداد ۱ تا ۶) و دیگری غیرمعمولی (که همیشه ۶ می‌آورد). به طور تصادفی یکی از تاس‌ها را انتخاب کرده و می‌اندازیم و عدد ۶ می‌آید. احتمال اینکه تاس غیرمعمولی را انتخاب کرده باشیم چقدر است؟

تعریف رویدادها:
\begin{itemize}
    \item رویداد A: انتخاب تاس غیرمعمولی
    \item رویداد B: آمدن عدد ۶
\end{itemize}
احتمالات داده شده:
\begin{itemize}
    \item $P(A) = \frac{1}{2}$ (احتمال یکسان برای انتخاب هر تاس)
    \item $P(B|A) = 1$ (تاس غیرمعمولی همیشه ۶ می‌آورد)
    \item $P(B|A^c) = \frac{1}{6}$ (تاس معمولی با احتمال ۱/۶ عدد ۶ می‌آورد)
\end{itemize}
اعمال قضیه بیز:
\[
P(B) = P(B|A)P(A) + P(B|A^c)P(A^c) = 1 \times \frac{1}{2} + \frac{1}{6} \times \frac{1}{2} = \frac{7}{12}
\]
\[
P(A|B) = \frac{P(B|A)P(A)}{P(B)} = \frac{1 \times 1/2}{7/12} = \frac{6}{7} \approx 85.7\%
\]
بنابراین، احتمال اینکه تاس انتخابی غیرمعمولی باشد، تقریباً ۸۵.۷ درصد است.

\section*{۳. دسته‌بندی‌کننده‌های بیز}
برای دسته‌بندی یک نمونه $d$، احتمال پسین قرار گرفتن در کلاس $c_j$ با فرمول زیر بیان می‌شود:
\[
P(c_j|d) = \frac{p(d|c_j)P(c_j)}{p(d)}
\]
که در آن:
\begin{itemize}
    \item احتمال پسین ($P(c_j|d)$): احتمال کلاس $c_j$ با فرض مشاهده نمونه $d$.
    \item درست‌نمایی ($p(d|c_j)$): احتمال مشاهده نمونه $d$ به شرط تعلق به کلاس $c_j$.
    \item احتمال پیشین ($P(c_j)$): احتمال وقوع کلاس $c_j$ در مجموعه‌داده (نرخ پایه).
    \item شواهد ($p(d)$): احتمال کل مشاهده نمونه $d$ در تمامی کلاس‌ها که به عنوان فاکتور نرمال‌سازی عمل می‌کند.
\end{itemize}

فرمول ساده‌شده برای دسته‌بندی (با حذف مخرج مشترک):
\[
P(c_j|d) \propto p(d|c_j) P(c_j)
\]
ما نمونه $d$ را به کلاسی با بیشترین احتمال پسین اختصاص می‌دهیم.

\section*{۴. نایو بیز با ویژگی‌های متعدد}
در صورتی که چندین ویژگی داشته باشیم، دسته‌بندی‌کننده‌های نایو بیز با فرض ساده‌لوحانه مستقل بودن توزیع ویژگی‌ها از یکدیگر، درست‌نمایی را به صورت حاصل‌ضرب احتمالات تک‌تک ویژگی‌ها تخمین می‌زنند:
\[
p(d|c_j) = p(d_1|c_j) \times p(d_2|c_j) \times \dots \times p(d_n|c_j)
\]

\subsection*{مثال میوه‌ها}
در یک مجموعه‌داده شامل ۱۰۰۰ میوه (موز، پرتقال، و سایر میوه‌ها)، می‌خواهیم میوه جدیدی با ویژگی‌های (بلند، شیرین، زرد) را دسته‌بندی کنیم. با محاسبه احتمالات پسین داریم:
\begin{itemize}
    \item $P(\text{Banana} | \text{Long, Sweet, Yellow}) = 0.8 \times 0.7 \times 0.9 \times 0.5 = 0.252$
    \item $P(\text{Orange} | \text{Long, Sweet, Yellow}) = 0 \times 0.5 \times 1 \times 0.3 = 0$
    \item $P(\text{Other} | \text{Long, Sweet, Yellow}) = 0.5 \times 0.75 \times 0.25 \times 0.2 = 0.01875$
\end{itemize}
از آنجا که احتمال پسین موز از همه بیشتر است، میوه جدید به عنوان موز دسته‌بندی می‌شود.

\section*{۵. مدل‌های زایشی}
مدلی که چگالی درست‌نمایی $P(\text{features}|L)$ را برای هر کلاس تعریف می‌کند، مدل زایشی (\lr{generative model}) نامیده می‌شود، زیرا فرآیند تصادفی فرضی تولید داده‌ها را مشخص می‌کند. در دسته‌بندی بیز، مشخص کردن این مدل زایشی گام اصلی آموزش است.

\section*{۶. نایو بیز گاوسی}
ساده‌ترین مدل نایو بیز، مدل نایو بیز گاوسی (\lr{Gaussian Naive Bayes}) است. در این مدل فرض می‌شود داده‌های مربوط به هر کلاس از یک توزیع گاوسی ساده بدون کوواریانس بین ابعاد استخراج شده‌اند. این مدل با محاسبه میانگین و انحراف معیار نقاط در هر کلاس برازش می‌شود.
مرز تصمیم‌گیری حاصل از مدل نایو بیز گاوسی عموماً درجه دوم (منحنی) است. این روش امکان دسته‌بندی احتمالی (\lr{probabilistic classification}) را نیز فراهم می‌کند.
\begin{figure}[h!]
  \centering
  \includegraphics[width=0.75\textwidth]{img_page_23_1.png}
  \caption{پیش‌بینی کلاس نقاط جدید بر روی فضای ویژگی}
\end{figure}




\section{۷. نایو بیز چندجمله‌ای}
یکی دیگر از مدل‌ها، نایو بیز چندجمله‌ای (\lr{Multinomial Naive Bayes}) است که در آن فرض می‌شود ویژگی‌ها از یک توزیع چندجمله‌ای ساده تولید شده‌اند. این توزیع تعداد رخدادها را در چند طبقه توصیف می‌کند و برای ویژگی‌هایی که نشان‌دهنده فراوانی یا تعداد شمارش هستند مناسب است.

کاربرد اصلی این روش در دسته‌بندی متون (\lr{text classification}) است، که در آن ویژگی‌ها مربوط به تعداد تکرار یا فرکانس کلمات در اسناد هستند.

\section*{۸. چه زمانی از نایو بیز استفاده کنیم؟}
با وجود فرضیات ساده‌لوحانه، دسته‌بندی‌کننده نایو بیز مزایای زیر را دارد:
\begin{itemize}
    \item آموزش و پیش‌بینی بسیار سریع.
    \item ارائه پیش‌بینی‌های احتمالی صریح.
    \item تفسیرپذیری بالا و پارامترهای قابل تنظیم بسیار کم.
\end{itemize}
این الگوریتم به ویژه برای داده‌های با ابعاد بسیار بالا (مانند متون) و دسته‌بندی‌های به خوبی تفکیک‌شده بسیار کارآمد عمل می‌کند.

\end{document}
